\section{Direct quartic strict vanishing in the second weight residue}
\label{sec:cubic-ramified-quartic-strict-vanishing}

The class certificate AC-G133 permits a second direct strict
vanishing theorem over the original field.  Here the minimal quartic
field is ramified at269, and the local unit calculation has absolute
ramification degree six.  The local parameter, both primes above269,
and both quartic projectors are computed explicitly.  The first
residue theorem AC-G130 is used only in the final combined count.
No automorphic endpoint, auxiliary host, native support, or model
condition is assumed in either strict vanishing statement below.

\paragraph{Characters and the unchanged strict condition.}
Put $p=269$, $K=\mathbb Q(\sqrt{-3})$, and let$t$ be the real cube root of$p$.
Fix a generator $\chi:G_K\to\mu_3$ for $K(t)/K$.
Let $\psi$ be the primitive quartic character modulo five with
$\psi(2)=i$, choosing the $p$-adic fourth root whose residue is82.
These fourth roots already belong to$\mathbb Q_p$ and to every
coefficient ring considered below.
Its arithmetic Galois avatar is $\Psi$; thus if
$g(\zeta_5)=\zeta_5^a$, then $\Psi(g)=\psi(a)$.
Write $c_p$ for the action on $p$-power roots of unity, so its
arithmetic Frobenius value at $q\ne p$ is$q$.
Let $\tau_p=\bar c_p^{134}$, regarded as a quadratic character.
It is the character of $\mathbb Q(\sqrt p)$, since $p\equiv1\pmod4$.

For $j=1,2$, $\sigma=1,-1$, and $e\ge1$ with
$e\equiv133\pmod{268}$, put
\[
       a_{j,\sigma,e}=\chi^{-j}\Psi^{-\sigma}c_p^{-e}.
 \tag{RQSV.1}
\]
For a finite discrete module, or for the discrete $E/O$ modules
below, use the literal strict group
\[
 \begin{split}
 \operatorname{Sel}_{\mathrm{str}}(K,V)=\ker\biggl(
 H^1_{\mathrm{cts}}(G_K,V)&\longrightarrow
           \prod_{v\mid p}H^1_{\mathrm{cts}}(G_{K_v},V)\\
       &\hspace{5mm}\times
           \prod_{v\nmid p,\infty}H^1_{\mathrm{cts}}(I_v,V)
                                  \biggr).
 \end{split}
 \tag{RQSV.2}
\]
There is zero restriction to full local cohomology at$p$, and
unramifiedness away from$p$.  A Bloch--Kato finite local condition
is not substituted for either of these requirements.

\begin{theorem}[AC-G134-residual: the ramified quartic residual sectors]
\label{thm:ACG134residual}
Let $k$ be any finite field of characteristic269 containing the
reduced character values.  For all the data above,
\[
          \operatorname{Sel}_{\mathrm{str}}
               (K,k(\bar a_{j,\sigma,e}))=0
              \qquad (j=1,2,\ \sigma=1,-1).
 \tag{RQSV.3}
\]
The class input is exactly AC-G133-quartic-class for
$L=\mathbb Q(t,w)$, where
$w^4+1345w^2+361805=0$, and $M=LK$.
Only its nontrivial cubic class-sector conclusion for$M$ is used.
The unit input is the single exact fifth global unit from that
certificate, with the independent local calculation below.
\end{theorem}

\begin{proof}
\proofstep{1. The actual minimal Kummer splitting field.}
Since $-134\equiv134\pmod{268}$, reduction of (RQSV.1) gives
\[
 \begin{split}
 \bar a_{j,\sigma,e}
      &=\bar\chi^{-j}\bar\Psi^{-\sigma}\tau_p\bar c_p,\\
 \eta_{j,\sigma}:=\bar a_{j,\sigma,e}/\bar c_p
      &=\bar\chi^{-j}(\bar\Psi\tau_p)^{-\sigma}.
 \end{split}
 \tag{RQSV.4}
\]
This uses the arithmetic Tate character, without replacing it by
its inverse.  The quartic factor is $\Psi\tau_p$, not$\Psi$ alone.

In the disjoint compositum $\mathbb Q(\zeta_5,\sqrt p)$ set
\[
       w=\sqrt p(\zeta_5-\zeta_5^{-1}),\qquad
       F=\mathbb Q(w),\quad L=\mathbb Q(t,w),\quad M=KL.
 \tag{RQSV.5}
\]
The element which squares $\zeta_5$ and fixes $\sqrt p$ has
quartic character value$i$; the element negating $\sqrt p$
has quadratic value$-1$.  Their common order-two kernel for
$\Psi\tau_p$ fixes$w$.  Thus $F$ is its cyclic quartic field,
with unique quadratic subfield$\mathbb Q(\sqrt5)$ and equation
\[
                        w^4+5pw^2+5p^2=0.
 \tag{RQSV.6}
\]
The cubic splitting field and$F$ have trivial intersection:
their only possible nontrivial abelian intersection is$K$,
which is ramified at three while$F$ is not.
Consequently
\[
        \operatorname{Gal}(M/\mathbb Q)=S_3\times C_4,
        \qquad \Delta=\operatorname{Gal}(M/K)=C_3\times C_4.
 \tag{RQSV.7}
\]
The product character in (RQSV.4) is faithful on$\Delta$.
Therefore$M$ is exactly its splitting field; adjoining$\mu_p$
is unnecessary.  In particular $\mu_p\not\subset M$, since
$[M:\mathbb Q]=24$ and $[\mathbb Q(\mu_p):\mathbb Q]=268$.

\proofstep{2. Inverse-character descent and all finite valuations.}
Let$S$ be the primes of$M$ above$p$, and define
\[
 \begin{split}
 V_{M,S}&=\{x\in M^\times:
       p\mid\operatorname{ord}_v(x)\text{ for }v\nmid p\}/M^{\times p},\\
 V_M&=\{x\in M^\times:
       p\mid\operatorname{ord}_v(x)\text{ for every finite }v\}/M^{\times p},\\
 \mathcal L_M&=\bigoplus_{v\mid p}M_v^\times/M_v^{\times p}.
 \end{split}
 \tag{RQSV.8}
\]
Prime-to-$p$ inflation--restriction and Kummer theory identify
the group in (RQSV.3) with
\[
 e_{\eta_{j,\sigma}^{-1}}\ker\left(
           V_{M,S}\otimes_{\mathbb F_p}k
             \longrightarrow\mathcal L_M\otimes_{\mathbb F_p}k
                            \right).
 \tag{RQSV.9}
\]
Indeed on$G_M$ the residual coefficient becomes$k(1)$, and the
natural action on$H^1(M,\mu_p)$ is twisted by$\eta_{j,\sigma}$.
Its invariant classes are the inverse-character part
\[
                \eta_{j,\sigma}^{-1}
                  =\bar\chi^j(\bar\Psi\tau_p)^\sigma.
 \tag{RQSV.10}
\]
The same argument applies to each local decomposition and inertia
group: all finite quotients involved have order dividing twelve.
It preserves the exact zero local classes.  For residue
characteristic different from$p$, Kummer unramifiedness is exactly
valuation divisible by$p$.  At$p$, a zero class is a local$p$th
power.  No single chosen completion replaces the full sum in
(RQSV.8).

A class in the kernel of (RQSV.9) also has valuation divisible
by$p$ at every place above$p$, and so lies in$V_M\otimes k$.
This remains true after coefficient extension, since$V_M$ is
the kernel of the remaining valuation map and field extension
is exact.  The equivariant ideal-factorization sequence is
\[
 0\longrightarrow E_M/E_M^p\longrightarrow V_M
   \longrightarrow\operatorname{Cl}(M)[p]\longrightarrow0,
 \qquad E_M=\mathcal O_M^\times .
 \tag{RQSV.11}
\]
For completeness, write$(x)=\mathfrak a^p$ and map$[x]$ to
$[\mathfrak a]$.  Its kernel is a unit modulo a$p$th power;
every$p$-torsion ideal class admits such a representative.
The cubic factor in (RQSV.10) is nontrivial, so AC-G133 annihilates
the relevant part of the last term in (RQSV.11).
Every strict class therefore comes from the corresponding
ordinary unit sector.  No assertion about an$S$-unit rank or a
class sector in a different field is used.

\proofstep{3. The exact unit multiplicities.}
The twelve complex places of$M$ are a free transitive$\Delta$-set,
after choosing representatives over a fixed complex embedding of$K$.
The equivariant logarithmic unit theorem gives the rational unit
character as the regular character of$\Delta$ minus the trivial
character.  These multiplicities remain valid modulo$p$:
after an unramified coefficient extension containing the twelfth
roots of unity, the character idempotents have denominators
dividing twelve and preserve the integral unit lattice.
Their ranks are unchanged on reduction.  The absence of$p$-power
roots of unity identifies this reduction with$E_M/E_M^p$.
No computed list is assumed to be a fundamental unit basis.

Let$b$ generate$C_4$ with$(\Psi\tau_p)(b)=i$ and fix$K,t$.
Let$\tau$ generate$C_3$ and satisfy$\tau(t)=\omega t$, where
$\omega^3=1$ and$\omega\ne1$.  In the following idempotents
$i$ means its residue$82\in\mathbb F_p$:
\[
 \begin{split}
 e_\sigma&=\frac14\sum_{a=0}^3 i^{-\sigma a}b^a,\\
 e_{\mathrm{nt}}&=1-\frac13(1+\tau+\tau^2),\\
 W_\sigma&=e_{\mathrm{nt}}e_\sigma(E_M/E_M^p),
                           \qquad \sigma=1,-1 .
 \end{split}
 \tag{RQSV.12}
\]
Each$W_\sigma$ has dimension two over$\mathbb F_p$.
Since$p\equiv2\pmod3$, the polynomial$X^2+X+1$ is irreducible
over$\mathbb F_p$, making$W_\sigma$ the simple nontrivial
two-dimensional$\mathbb F_p[C_3]$-module.
After extension to$k$, its two cubic-character lines are exactly
the two characters in (RQSV.10).

\proofstep{4. The ramified local parameter and both completions.}
Put
\[
             R=\frac{w^2}{p},\qquad u=-\frac wt .
 \tag{RQSV.13}
\]
The exact identities are
\[
 \begin{gathered}
 R^2+5R+5=0,\qquad
 u^6=pR^3,\qquad t=u^2/R,\qquad w=-u^3/R,\\
 b(R)=-5-R,\qquad b(w)=(-3-R)w,\qquad b^2(w)=-w.
 \end{gathered}
 \tag{RQSV.14}
\]
For example$\zeta_5+\zeta_5^{-1}=-3-R$ in the embedding
used in (RQSV.5).  Its product with its quadratic conjugate is$-1$,
which also checks the asserted square of$b$.

Modulo$p$ the polynomial$X^2+5X+5$ has the two simple roots
\[
                               69,\qquad195.
 \tag{RQSV.15}
\]
Hensel lifts place$R$ in$\mathbb Q_p$ in either embedding.
The equation$u^6=pR^3$ is Eisenstein there.  Also
$K_p/\mathbb Q_p$ is unramified quadratic.  It follows that$M$
has exactly two primes over$p$, each with absolute ramification
degree six and residue degree two.  The two roots in (RQSV.15)
label those primes.  In each completion$u$ is a uniformizer and
$\tau(u)=\omega^{-1}u$.
These statements differ from the absolute degree-three calculation
in AC-G130; no unramified quartic field is substituted here.

\proofstep{5. One exact global unit and its local expansions.}
Use the fifth unit$U$ of AC-G133.  To fix it without numerical
embeddings, write
\[
                U=\frac1{6p^2}
                    \sum_{a=0}^2\sum_{b=0}^3 N_{ab}t^a w^b,
 \tag{RQSV.16}
\]
with the following integer numerators:
\[
\begin{array}{cc|r}
a&b&N_{ab}\\ \hline
0&0&-21188087364390629115228010464\\
0&1&-1679460225000426189781561242\\
0&2&-56995709003186289184103540\\
0&3&-4517728189718195097835902\\
1&0&-5387894948998345378336713084\\
1&1&221199065653505753324695437\\
1&2&-14493373250324747359323571\\
1&3&595022910412010111333019\\
2&0&1343104928559150944288897379\\
2&1&6036574145212979063569527\\
2&2&3612936414800857297005739\\
2&3&16238308009682117277954
\end{array}
 \tag{RQSV.17}
\]
The exact rational multiplication and raw-generator verifiers
prove a monic integral degree-twelve relation with constant one
for this very element.  Hence both$U$ and$U^{-1}$ are integral.
This global-unit fact is independent of the provisional PARI
class output or completeness of its candidate units.

Set$B=-3-R$ and$\lambda_a=(1,B,-1,-B)_a$ for$0\le a<4$.
To expand$b^a(U)$, replace$w$ by$\lambda_a w$ in (RQSV.16),
then use (RQSV.14).  A monomial$t^i w^j$ becomes
$(-1)^j u^{2i+3j}/R^{i+j}$; reduce$u^6$ to$pR^3$.
This gives six coefficients in$\mathbb Q[R]/(R^2+5R+5)$.
The verifier reduces every such coefficient and checks that its
two rational denominators are prime to$p$, before reducing at
either root in (RQSV.15).
In particular the constant and linear coefficients are
\[
 \begin{split}
 A_a&=\frac{N_{00}+N_{02}\lambda_a^2pR}{6p^2},\\
 C_a&=-\frac{N_{21}\lambda_ap+N_{23}\lambda_a^3p^2R}{6p^2}.
 \end{split}
 \tag{RQSV.18}
\]
The linear formula retains the contributions from both$t^2w$
and$t^2w^3$.  Dropping the former would give incorrect numbers.
Exact reduction gives$A_a=1$ modulo$u$ for every$a$, and
\[
 \begin{array}{c|cc}
 a&C_a\text{ at }R\equiv69&C_a\text{ at }R\equiv195\\ \hline
 0&176&48\\
 1&41&122\\
 2&93&221\\
 3&228&147
 \end{array}
 \qquad b^a(U)=1+C_a u+O(u^2).
 \tag{RQSV.19}
\]
All displayed residues lie in the prime field of the degree-two
residue field; the chosen cubic root of unity lies in that latter
field.

Represent the two idempotents in (RQSV.12) by the integer weights
\[
 (m_{+,a})=(202,114,67,155),\qquad
 (m_{-,a})=(202,155,67,114).
 \tag{RQSV.20}
\]
Form the actual global units
\[
 A_\sigma=\prod_{a=0}^3 b^a(U)^{m_{\sigma,a}},\qquad
 V_\sigma=\frac{A_\sigma^3}
                {A_\sigma\tau(A_\sigma)\tau^2(A_\sigma)}.
 \tag{RQSV.21}
\]
Their classes satisfy$[V_\sigma]=3e_{\mathrm{nt}}e_\sigma[U]$
and belong to$W_\sigma$.  Only integer powers occur in these
global units; the denominators three and four belong to the
prime-to-$p$ character algebra.
The cubic norm in (RQSV.21) has constant one and zero linear term,
because$1+\omega+\omega^2=0$ and$\tau(u)=\omega^{-1}u$.
Thus the exact expansions are
\[
 \begin{array}{c|cc}
       &R\equiv69&R\equiv195\\ \hline
 V_+&1+63u+O(u^2)&1+130u+O(u^2)\\
 V_-&1+196u+O(u^2)&1+14u+O(u^2).
 \end{array}
 \tag{RQSV.22}
\]
For example each coefficient is$3\sum_a m_{\sigma,a}C_a$ modulo$p$.
The verifier also multiplies the actual products and norm in the
dual-number ring over$\mathbb F_p[s]/(s^2+3)$ with
$\omega=(-1+s)/2$, independently checking that formula.

None of the four local units in (RQSV.22) is a$p$th power.
If$V=y^p$ reduces to one, then$y$ also reduces to one, since
raising to$p$ is bijective on the residue multiplicative group.
Write$y=1+x$ with$v_u(x)\ge1$.
For$1\le j<p$, the coefficient$\binom pj$ has$u$-valuation six.
Therefore
\[
        v_u((1+x)^p-1)\ge
                  \min(6+v_u(x),p\,v_u(x))\ge7.
 \tag{RQSV.23}
\]
This contradicts a nonzero linear term in (RQSV.22).

\proofstep{6. Injectivity and the residual conclusion.}
The localization map from$E_M/E_M^p$ to the sum over both
$p$-adic completions is$\Delta$-equivariant.
Equation (RQSV.22) makes its restriction to each$W_\sigma$
nonzero.  By its simplicity as an$\mathbb F_p[C_3]$-module,
that restriction is injective.  Coefficient extension preserves
injectivity, so both cubic-character lines over$k$ inject.
Steps2--3 placed every strict residual class in those lines.
Its local image is zero by definition, so it is zero.
This proves (RQSV.3) for both$j$ and both$\sigma$.
\end{proof}

\begin{theorem}[AC-G134-strict: original coefficients and the second lift family]
\label{thm:ACG134strict}
Let$E/\mathbb Q_p$ be any finite extension with ring$O$,
uniformizer$\varpi$, and residue field$k$, containing the values
of$\chi$.  Keep this original$O$ throughout.
Let$\omega_p$ be the Teichmuller lift of$\bar c_p$ and let
$\kappa:G_K\to1+\mathfrak m_O$ be any continuous character
unramified at every finite place away from$p$.
For$r\in\mathbb Z$, $e\ge1$ with$e+r\equiv133\pmod{268}$, put
\[
 \begin{split}
 \psi_{\sigma,r,\kappa}&=\Psi^\sigma\omega_p^r\kappa,\\
 D_{j,\sigma,r,\kappa,e}
       &=(E/O)(\chi^{-j}\psi_{\sigma,r,\kappa}^{-1}c_p^{-e}).
 \end{split}
 \tag{RQSV.24}
\]
Then
\[
      \operatorname{Sel}_{\mathrm{str}}
             (K,D_{j,\sigma,r,\kappa,e})=0
                  \qquad(j=1,2,\ \sigma=1,-1).
 \tag{RQSV.25}
\]
Both its$O$-length and its length divided by
$e_O=\operatorname{length}_O(O/pO)$ are zero.
No condition$e_O=1$, finite order of$\kappa$, or enlargement of
the original coefficient ring is imposed.
\end{theorem}

\begin{proof}
Write$D$ for the module and$a$ for its character.
The residual character is
$\bar\chi^{-j}\bar\Psi^{-\sigma}\bar c_p^{-(e+r)}$,
whose strict group is zero by the preceding theorem.
Multiplication by the original uniformizer gives
\[
 0\longrightarrow k(\bar a)\longrightarrow D
       \xrightarrow{\varpi}D\longrightarrow0,
 \qquad \bar x\longmapsto x/\varpi+O.
 \tag{RQSV.26}
\]
This is a uniformizer kernel, not the$p$-torsion kernel over
a potentially ramified coefficient ring.

At$p$, the inertia image of$\chi$ has order three, whereas
that of$\bar c_p$ has order dividing268 and$\Psi$ is unramified.
The reduction of$\kappa$ is one.  Raising an inertia element
with nontrivial cubic value to the268th power kills the rational
cyclotomic factor and preserves the nontrivial cubic value.
Some inertia element therefore has$\bar a(g)\ne1$, making
$a(g)-1$ a unit of$O$.  Hence
\[
                        D^{G_{K_p}}=D^{G_K}=0.
 \tag{RQSV.27}
\]
At every finite$v\nmid p$, the factors$c_p,\omega_p,\kappa$
are unramified.  The remaining inertia character has finite
order dividing twelve.  Its invariant module is$D$ if the
character is trivial, and zero otherwise: a nontrivial
prime-to-$p$ root of unity differs from one by a unit.
Thus every local invariant module in (RQSV.2) is
$\varpi$-divisible, including the places over three and five.

The global long exact sequence of (RQSV.26), using (RQSV.27),
identifies$H^1(G_K,k(\bar a))$ with$H^1(G_K,D)[\varpi]$.
For each local group$J$ in (RQSV.2), divisibility of$D^J$
makes$H^1(J,k(\bar a))\to H^1(J,D)$ injective.
The lift of a strict uniformizer-torsion class is consequently
a residual strict class, and is zero.
Explicitly, a local coboundary$dx$ with values in$D[\varpi]$
has$\varpi x\in D^J$; choose$x_0\in D^J$ with
$\varpi x_0=\varpi x$.  Then$x-x_0\in D[\varpi]$ gives the
same coboundary, which verifies the literal local-zero condition.

Finally every continuous cocycle from compact$G_K$ to discrete
$E/O$ has finite image.  A power of$\varpi$ kills every value,
so every cohomology class is uniformizer-primary torsion.
A nonzero strict class would have a last nonzero uniformizer
multiple, giving nonzero strict$\varpi$-torsion.  This has just
been excluded, proving (RQSV.25) without a finiteness or Iwasawa
assumption.
\end{proof}

\paragraph{Both residues and a precise finite character count.}
Combining AC-G130 with (RQSV.25) permits either
\[
                         e+r\equiv133\ \text{or}\ 267\pmod{268}.
 \tag{RQSV.28}
\]
For fixed$e\ge1$ and$m\ge1$, there are two distinct tame
exponents$r$ in (RQSV.28), and$p^{m-1}$ wild characters of
conductor dividing$p^m$.  Each of the two quartic labels at
five therefore supplies$2p^{m-1}$ characters.  Altogether
\[
                              4\cdot269^{m-1}
 \tag{RQSV.29}
\]
distinct sufficient Dirichlet labels of conductor dividing
$5\cdot269^m$ have strict vanishing for both cubic labels.
The exponents differing by134 give distinct characters at$p$;
the two quartic labels differ at five.  Restriction to$G_K$
does not identify any two: a quotient trivial there is either
one or$\eta_K$, and$\eta_K$ is ramified at three whereas all
these quotients are unramified there.
Every label has parity
$\psi^\sigma(-1)(-1)^r=(-1)^{r+1}=(-1)^e$,
since both permitted values of$e+r$ are odd and a wild
odd-prime-order character is trivial on$-1$.
This counts sufficient character labels.  It is not an
arithmetic family-average or a numerical window bound.

\paragraph{Finite evidence, retained sources, and geometry.}
The class input is the reviewed AC-G133 analytic interval,
exact units,269-saturation and reflection descent.  No
provisional PARI class output is used as a class certificate.
The local script\texttt{verify\_ramified\_local\_projectors.py}
first works over$\mathbb Q[R]/(R^2+5R+5)$, checks all six
$p$-integral coefficients, then verifies the two projector
products at both residue roots in dual numbers.
Its exact output is
\texttt{ramified\_local\_projectors\_certificate.json}.
The global unit is independently fixed by (RQSV.16)--(RQSV.17).
Only its monic integral constant-one annihilator is needed;
there is no assumption of a complete unit basis.

The retained Kummer exact sequence is Stacks Project tag03PK,
Lemma59.28.1 and the ensuing unit/Picard interface.  The
inverse-character descent, all-primes ideal sequence, unit
multiplicity, local calculation and original-uniformizer
argument have been proved above.  Dirichlet's equivariant
logarithmic unit theorem is retained at its stated scope.
Geometrically, (RQSV.9) describes the indicated character
parts of Kummer torsors over$\operatorname{Spec}\mathcal O_M[1/p]$
which are trivial at both$p$-adic completions.  The class-sector
certificate removes their ideal obstruction; the explicit
unit map is injective on both cubic-character lines.
The conclusion concerns these actual torsors and their
original-coefficient lifts.  It supplies no new automorphic
support, mixed Bloch--Kato vanishing, class certificate for a
different reflection field, or H04.4 arithmetic window inequality.
